% Options for packages loaded elsewhere
\PassOptionsToPackage{unicode}{hyperref}
\PassOptionsToPackage{hyphens}{url}
%
\documentclass[
  10pt,
]{article}
\usepackage{amsmath,amssymb}
\usepackage{iftex}
\ifPDFTeX
  \usepackage[T1]{fontenc}
  \usepackage[utf8]{inputenc}
  \usepackage{textcomp} % provide euro and other symbols
\else % if luatex or xetex
  \usepackage{unicode-math} % this also loads fontspec
  \defaultfontfeatures{Scale=MatchLowercase}
  \defaultfontfeatures[\rmfamily]{Ligatures=TeX,Scale=1}
\fi
\usepackage{lmodern}
\ifPDFTeX\else
  % xetex/luatex font selection
\fi
% Use upquote if available, for straight quotes in verbatim environments
\IfFileExists{upquote.sty}{\usepackage{upquote}}{}
\IfFileExists{microtype.sty}{% use microtype if available
  \usepackage[]{microtype}
  \UseMicrotypeSet[protrusion]{basicmath} % disable protrusion for tt fonts
}{}
\makeatletter
\@ifundefined{KOMAClassName}{% if non-KOMA class
  \IfFileExists{parskip.sty}{%
    \usepackage{parskip}
  }{% else
    \setlength{\parindent}{0pt}
    \setlength{\parskip}{6pt plus 2pt minus 1pt}}
}{% if KOMA class
  \KOMAoptions{parskip=half}}
\makeatother
\usepackage{xcolor}
\usepackage{color}
\usepackage{fancyvrb}
\newcommand{\VerbBar}{|}
\newcommand{\VERB}{\Verb[commandchars=\\\{\}]}
\DefineVerbatimEnvironment{Highlighting}{Verbatim}{commandchars=\\\{\}}
% Add ',fontsize=\small' for more characters per line
\newenvironment{Shaded}{}{}
\newcommand{\AlertTok}[1]{\textcolor[rgb]{1.00,0.00,0.00}{\textbf{#1}}}
\newcommand{\AnnotationTok}[1]{\textcolor[rgb]{0.38,0.63,0.69}{\textbf{\textit{#1}}}}
\newcommand{\AttributeTok}[1]{\textcolor[rgb]{0.49,0.56,0.16}{#1}}
\newcommand{\BaseNTok}[1]{\textcolor[rgb]{0.25,0.63,0.44}{#1}}
\newcommand{\BuiltInTok}[1]{\textcolor[rgb]{0.00,0.50,0.00}{#1}}
\newcommand{\CharTok}[1]{\textcolor[rgb]{0.25,0.44,0.63}{#1}}
\newcommand{\CommentTok}[1]{\textcolor[rgb]{0.38,0.63,0.69}{\textit{#1}}}
\newcommand{\CommentVarTok}[1]{\textcolor[rgb]{0.38,0.63,0.69}{\textbf{\textit{#1}}}}
\newcommand{\ConstantTok}[1]{\textcolor[rgb]{0.53,0.00,0.00}{#1}}
\newcommand{\ControlFlowTok}[1]{\textcolor[rgb]{0.00,0.44,0.13}{\textbf{#1}}}
\newcommand{\DataTypeTok}[1]{\textcolor[rgb]{0.56,0.13,0.00}{#1}}
\newcommand{\DecValTok}[1]{\textcolor[rgb]{0.25,0.63,0.44}{#1}}
\newcommand{\DocumentationTok}[1]{\textcolor[rgb]{0.73,0.13,0.13}{\textit{#1}}}
\newcommand{\ErrorTok}[1]{\textcolor[rgb]{1.00,0.00,0.00}{\textbf{#1}}}
\newcommand{\ExtensionTok}[1]{#1}
\newcommand{\FloatTok}[1]{\textcolor[rgb]{0.25,0.63,0.44}{#1}}
\newcommand{\FunctionTok}[1]{\textcolor[rgb]{0.02,0.16,0.49}{#1}}
\newcommand{\ImportTok}[1]{\textcolor[rgb]{0.00,0.50,0.00}{\textbf{#1}}}
\newcommand{\InformationTok}[1]{\textcolor[rgb]{0.38,0.63,0.69}{\textbf{\textit{#1}}}}
\newcommand{\KeywordTok}[1]{\textcolor[rgb]{0.00,0.44,0.13}{\textbf{#1}}}
\newcommand{\NormalTok}[1]{#1}
\newcommand{\OperatorTok}[1]{\textcolor[rgb]{0.40,0.40,0.40}{#1}}
\newcommand{\OtherTok}[1]{\textcolor[rgb]{0.00,0.44,0.13}{#1}}
\newcommand{\PreprocessorTok}[1]{\textcolor[rgb]{0.74,0.48,0.00}{#1}}
\newcommand{\RegionMarkerTok}[1]{#1}
\newcommand{\SpecialCharTok}[1]{\textcolor[rgb]{0.25,0.44,0.63}{#1}}
\newcommand{\SpecialStringTok}[1]{\textcolor[rgb]{0.73,0.40,0.53}{#1}}
\newcommand{\StringTok}[1]{\textcolor[rgb]{0.25,0.44,0.63}{#1}}
\newcommand{\VariableTok}[1]{\textcolor[rgb]{0.10,0.09,0.49}{#1}}
\newcommand{\VerbatimStringTok}[1]{\textcolor[rgb]{0.25,0.44,0.63}{#1}}
\newcommand{\WarningTok}[1]{\textcolor[rgb]{0.38,0.63,0.69}{\textbf{\textit{#1}}}}
\usepackage{longtable,booktabs,array}
\usepackage{calc} % for calculating minipage widths
% Correct order of tables after \paragraph or \subparagraph
\usepackage{etoolbox}
\makeatletter
\patchcmd\longtable{\par}{\if@noskipsec\mbox{}\fi\par}{}{}
\makeatother
% Allow footnotes in longtable head/foot
\IfFileExists{footnotehyper.sty}{\usepackage{footnotehyper}}{\usepackage{footnote}}
\makesavenoteenv{longtable}
\setlength{\emergencystretch}{3em} % prevent overfull lines
\providecommand{\tightlist}{%
  \setlength{\itemsep}{0pt}\setlength{\parskip}{0pt}}
\setcounter{secnumdepth}{-\maxdimen} % remove section numbering
\usepackage{fontspec}
\setmainfont{DejaVu Serif}[Scale=0.88]
\setsansfont{DejaVu Sans}[Scale=0.88]
\setmonofont{DejaVu Sans Mono}[Scale=0.8]
\usepackage{unicode-math}
\setmathfont{latinmodern-math.otf}
\usepackage[margin=2.2cm]{geometry}
\usepackage{longtable,booktabs}
\setlength{\emergencystretch}{3em}
\sloppy
\usepackage{newunicodechar}
\newunicodechar{∎}{\ensuremath{\blacksquare}}
\newunicodechar{∖}{\ensuremath{\backslash}}
\newunicodechar{≅}{\ensuremath{\cong}}
\newunicodechar{⋯}{\ensuremath{\cdots}}
\newunicodechar{𝒬}{\ensuremath{\mathcal{Q}}}
\ifLuaTeX
  \usepackage{selnolig}  % disable illegal ligatures
\fi
\IfFileExists{bookmark.sty}{\usepackage{bookmark}}{\usepackage{hyperref}}
\IfFileExists{xurl.sty}{\usepackage{xurl}}{} % add URL line breaks if available
\urlstyle{same}
\hypersetup{
  hidelinks,
  pdfcreator={LaTeX via pandoc}}

\title{Uniform families in the shape of the D = 8���10 Haar records}
\author{}
\date{2026-10-27 (draft)}

\begin{document}
\maketitle

\hypertarget{uniform-families-in-the-shape-of-the-d-810-haar-records-the-lps-family-a-43-obstruction-and-what-a-cyclic-factor-does}{%
\section{Uniform families in the shape of the D = 8--10 Haar records:
the LPS family, a 4/3 obstruction, and what a cyclic factor
does}\label{uniform-families-in-the-shape-of-the-d-810-haar-records-the-lps-family-a-43-obstruction-and-what-a-cyclic-factor-does}}

Doğa Özden's degree--diameter lab, with Claude Code (job
\texttt{thm-haar-families}, 2026-10-27). Written for Gabriela
Araujo-Pardo and Ignacio López. Draft; corrections welcome.

\textbf{Labels.} Every statement carries one of three labels.

\begin{itemize}
\tightlist
\item
  \textbf{Theorem, Proposition, Lemma, Corollary}: complete proof in
  this note, or a cited result named in place (marked \emph{cited}, with
  the source and how we know it: \textbf{{[}read{]}} primary source
  read, \textbf{{[}abstract{]}} abstract only, \textbf{{[}secondary{]}}
  known from a book, survey or citation).
\item
  \textbf{Computer-checked}: a statement about finitely many explicit
  graphs, verified by programs that are named, with logs.
\item
  \textbf{Heuristic}: a model, a measurement or an interpretation. Not
  used in any proof.
\end{itemize}

\textbf{The question.} The lab's records at D = 8, 9, 10
(\texttt{research/theorems/HAAR\_D8\_10.md} on branch
\texttt{claude/thm-haar-d8-10-sf7f0a}, cited below as {[}H{]}) are Haar
graphs of SL(2, q) × \(Z_{m}\) or PSL(2, q) × \(Z_{m}\) whose connection
set S is closed under an automorphism of order 3, each found by computer
search for one group. Can they be turned into infinite families, with S
given by one formula in q (and m) and a proved diameter?

\textbf{The answer, stated plainly.}

\begin{enumerate}
\def\labelenumi{\arabic{enumi}.}
\tightlist
\item
  \emph{There is an explicit infinite family of exactly this shape},
  with m = 1: the connection sets of Lubotzky, Phillips and Sarnak
  (LPS). For every odd prime p and every prime q \textgreater{} p, the p
  + 1 LPS matrices in PGL(2, q) are permuted by conjugation with an
  element of order 3 of PSL(2, q), with cycle type 1+1+3+⋯+3 or 3+⋯+3:
  the cycle types that the orbit search used at Δ = 6, 8, 12, 14
  (Theorem 2.6). Their diameter is at most 2 log\_p n + O(1)
  (\emph{cited}, LPS).
\item
  \emph{This family provably cannot reach the records.} The vertex ψ(i)
  is at distance exactly the least even L with 2p\(^{L/2}\)
  \textgreater{} q² (Theorem 3.2: lower bound for every odd p, equality
  for p ≡ 1 mod 4), so the diameter is at least (4/3) log\_p n − 2
  log\_p 2 (Corollary 3.3, the bound of N. T. Sardari, here with a
  two-paragraph proof), and a member of diameter ≤ D has at most
  √2·p\(^{1−D/4}\) of the bipartite Moore bound (Corollary 3.4).
  Computed: the LPS graphs are 2.5 to 19 times below the 2013 table at D
  = 8--10 (Appendix B), and at odd D the largest member is usually
  exactly the one that the vertex ψ(i) allows (8 of 13 cells; the next
  prime below in 3 more).
\item
  \emph{A cyclic factor removes this obstruction in practice, but not as
  a formula.} The obstruction survives every cover (Proposition 4.1),
  but it bounds q, not m. LPS × \(Z_{m}\) with labels chosen for each
  (q, m) by a squad reaches 99.4 \% of the table at (6,10) and 98.9 \%
  at (8,8) (Computer-checked 4.2), still below the lab's records.
\item
  \emph{Other uniform formulas behave like one random draw.} Integer
  lifts, primitive-root and Dickson formulas, the conic model, torus
  orbits, and the records' own sets moved to other q (18 squad cards,
  §5) give an excess diam − \(log_{Δ−1}\)\textbar G\textbar{} within
  about 0.3 of a single random draw at the same Δ (about 2.9--3.4 at Δ =
  6--8), against about 2.0--2.3 for the records. Heuristically, the
  records' advantage is selection among millions of sets, which no
  single formula provides (§5--6).
\item
  \emph{What would be needed is open.} A uniform family of record
  quality at fixed Δ would be an explicit family of bounded-degree
  graphs with diameter (1 + o(1)) \(log_{Δ−1}\) n.~None is known; for
  the LPS graphs Sardari proved that it fails (§6).
\end{enumerate}

No infinite family above the table is proved in this note.

\hypertarget{contents}{%
\subsection{Contents}\label{contents}}

\begin{enumerate}
\def\labelenumi{\arabic{enumi}.}
\setcounter{enumi}{-1}
\tightlist
\item
  Summary of statements
\item
  Setting, notation and projective coordinates
\item
  The LPS family has the records' shape
\item
  The 4/3 obstruction
\item
  Cyclic factors
\item
  Other uniform families (squad round 1)
\item
  Why no record-quality uniform family is in sight
\item
  Prior art
\item
  How to re-check
\end{enumerate}

Appendix A. The squad rounds. Appendix B. The LPS family cell by cell.

\begin{center}\rule{0.5\linewidth}{0.5pt}\end{center}

\hypertarget{summary-of-statements}{%
\subsection{0. Summary of statements}\label{summary-of-statements}}

\begin{longtable}[]{@{}
  >{\raggedright\arraybackslash}p{(\columnwidth - 4\tabcolsep) * \real{0.3333}}
  >{\raggedright\arraybackslash}p{(\columnwidth - 4\tabcolsep) * \real{0.3333}}
  >{\raggedright\arraybackslash}p{(\columnwidth - 4\tabcolsep) * \real{0.3333}}@{}}
\toprule\noalign{}
\begin{minipage}[b]{\linewidth}\raggedright
\#
\end{minipage} & \begin{minipage}[b]{\linewidth}\raggedright
statement
\end{minipage} & \begin{minipage}[b]{\linewidth}\raggedright
label
\end{minipage} \\
\midrule\noalign{}
\endhead
\bottomrule\noalign{}
\endlastfoot
Prop. 1.1 & at fixed (Δ, D) only finitely many groups; for Q × \(Z_{m}\)
with Q perfect, m ≤ A(Δ, D) & Proposition (from {[}H{]}) \\
Lemma 2.3 & ψ: H(Z) → \(M_{2}\)(\(F_{q}\)) is a ring homomorphism with
kernel qH(Z); det ψ(α) = N(α); ψ(ᾱ) = adj ψ(α) & Lemma \\
Lemma 2.4 & for q \textgreater{} p the p + 1 matrices ψ(α), α ∈
\(S_{p}\), are distinct in PGL(2, q), and ψ(\(S_{p}\)) is closed under
inversion & Lemma \\
Lemma 2.5 & structure of the component \(Γ_{p,q}\); diam \(Γ_{p,q}\) =
ecc(I, A) & Lemma \\
Thm. 2.6 & conjugation by h = ψ(1+i+j+k), of order 3 in PSL(2, q),
permutes ψ(\(S_{p}\)) with cycle type 1+1+3\(^{(p−1)/3}\) (p ≡ 1 mod 3),
3\(^{(p+1)/3}\) (p ≡ 2 mod 3), 1+3 (p = 3) & Theorem \\
Thm. 2.7 & \(Γ_{p,q}\) is connected and Ramanujan; diam ≤ 2 log\_p n +
O(1) & \emph{cited} (LPS) \\
Lemma 3.1 & an even walk from (I, A) ends at ψ(β) with N(β) = p\(^{L}\)
and β ≡ 1 mod 2H(Z) & Lemma \\
Thm. 3.2 & d((I,A), (ψ(i),A)) ≥ L₀ := min\{L even : 2p\(^{L/2}\)
\textgreater{} q²\}, with equality if p ≡ 1 mod 4 & Theorem (equality
uses a cited factorization) \\
Cor. 3.3 & diam \(Γ_{p,q}\) \textgreater{} (4/3) log\_p n − 2 log\_p 2 &
Corollary \\
Cor. 3.4 & if diam \(Γ_{p,q}\) ≤ D then n \textless{} 2\(^{3/2}\)
p\(^{3D/4}\) and n \textless{} √2·p\(^{1−D/4}\)·M\(^{b}\)(p+1, D) &
Corollary \\
3.5 & equality in Thm. 3.2 in all 251 cases p ∈ \{3, 5, 7, 11, 13, 17,
19, 29\}, p \textless{} q ≤ 160 (including p ≡ 3 mod 4); ecc − L₀ ∈ \{0,
1, 2, 3\} & Computer-checked \\
App. B & the LPS family cell by cell; at odd D its largest member is the
one Thm. 3.2 allows in 8 of 13 cells & Computer-checked \\
Prop. 3.6 & the same lower bound in any definite quaternion order with
an orthogonal unit γ & Proposition \\
Prop. 4.1 & the bound of Thm. 3.2 holds in every cover H(PGL(2,q) × K,
S̃) of the LPS graph & Proposition \\
4.2 & LPS × \(Z_{m}\) with searched labels: 99.4 \% of the table at
(6,10), 98.9 \% at (8,8) & Computer-checked \\
§5 & uniform formulas of six kinds: excess about that of one random draw
& Computer-checked (scores), Heuristic (reading) \\
\end{longtable}

\begin{center}\rule{0.5\linewidth}{0.5pt}\end{center}

\hypertarget{setting-notation-and-projective-coordinates}{%
\subsection{1. Setting, notation and projective
coordinates}\label{setting-notation-and-projective-coordinates}}

\hypertarget{haar-graphs}{%
\subsubsection{1.1 Haar graphs}\label{haar-graphs}}

We use the conventions of {[}H{]} §1--2. For a finite group G and S ⊆ G
with \textbar S\textbar{} = Δ, the Haar graph H(G, S) has vertex set G ×
\{A, B\}, and (g, A) is adjacent to (sg, B) for every s ∈ S. It is
bipartite and Δ-regular, the right translations \(ρ_{h}\)(g, X) = (gh,
X) are automorphisms acting regularly on each side ({[}H{]} Lemma 1.2),
and its diameter is max(ecc(e, A), ecc(e, B)) ({[}H{]} Cor. 2.3). A walk
of length L from (e, A) uses letters \(t_{1}\), \ldots, \(t_{L}\) ∈ S
and ends at (w, A) or (w, B) with

\[w = \cdots t_4^{-1}\, t_3\, t_2^{-1}\, t_1 \qquad (\text{an } A\text{-word}),\]

and the walk backtracks exactly when two consecutive letters are equal
({[}H{]} Prop. 2.2). Lemma 3.1 of {[}H{]}: H(G, S) is isomorphic to H(G,
aψ(S)\(^{±1}\)b) for a, b ∈ G and ψ ∈ Aut(G); the map (g, A) ↦ (g, B),
(h, B) ↦ (h, A) is an isomorphism H(G, S) → H(G, S⁻¹).

Throughout, M\(^{b}\)(Δ, D) = 2 \(Σ_{j=0}\)\(^{D−1}\) (Δ − 1)\(^{j}\) is
the bipartite Moore bound, and n denotes the number of vertices of the
graph at hand.

\hypertarget{projective-coordinates}{%
\subsubsection{1.2 Projective
coordinates}\label{projective-coordinates}}

Let q be an odd prime. An element of PGL(2, q) is a nonsingular 2 × 2
matrix up to a nonzero scalar, that is, a point

\[[X] = [a : b : c : d] \in \mathrm{PG}(3, q), \qquad X = \begin{pmatrix} a & b \\ c & d\end{pmatrix},\]

not on the quadric 𝒬 : ad − bc = 0. The quadric 𝒬 is a hyperbolic
quadric Q⁺(3, q) with (q + 1)² points, so PGL(2, q) is the set of the (q
+ 1)(q² + 1) − (q + 1)² = q(q² − 1) points of PG(3, q) off 𝒬. Whether ad
− bc is a square does not depend on the representative, and PSL(2, q) is
the set of the q(q² − 1)/2 points with ad − bc a nonzero square. For M ∈
GL(2, q), left multiplication {[}X{]} ↦ {[}MX{]} is a collineation of
PG(3, q) that fixes 𝒬 (det MX = det M det X).

So a Haar graph H(PGL(2, q), S) is the bipartite graph on two copies of
PG(3, q) ∖ 𝒬 in which \textbf{the point {[}X{]} of the first copy is
joined to the points {[}MX{]}, M ∈ S, of the second copy}; and the point
{[}Y{]} of the second copy is joined to the points {[}adj(M)·Y{]}, M ∈
S, of the first, where adj {[}{[}a, b{]}, {[}c, d{]}{]} = {[}{[}d,
−b{]}, {[}−c, a{]}{]} is the inverse up to a scalar. These are the exact
neighbourhoods of every vertex, once S is written down (§2, Example
2.8).

\hypertarget{the-records-shape}{%
\subsubsection{1.3 The records' shape}\label{the-records-shape}}

{[}H{]} Construction 3.2: S is closed under an automorphism α of order
k, here k = 3, and splits into f fixed letters and 3-cycles; for G = Q ×
\(Z_{m}\) the automorphism is α(x, z) = (gxg⁻¹, rz) with g of order 3
modulo the centre of Q and r of multiplicative order 3 mod m. The cycle
types of the records are in {[}H{]} §5.1: 3+3 at (6,9), (6,10); 1+3+3 at
Δ = 7; 1+1+3+3 at (8,8), (8,9), (8,10); 3+3+3 at (9,8), (9,9), (9,10);
3+3+3+3 at (12,9), (12,10); 1+1+3+3+3+3 at (14,9), (14,10); and so on.
Since \textbar S\textbar{} = f + 3c, f ≡ Δ (mod 3); every record with k
= 3 uses the least such f ∈ \{0, 1, 2\}.

By a \textbf{uniform family} we mean, informally, one rule that produces
S from q (and m) without a search for each q: a formula, as for the
Singer difference sets at D = 3. All the families below are uniform in
this sense, except where a search over labels is stated.

\hypertarget{where-a-uniform-family-can-live}{%
\subsubsection{1.4 Where a uniform family can
live}\label{where-a-uniform-family-can-live}}

\textbf{Proposition 1.1.} (a) For fixed Δ ≥ 2 and D, only finitely many
groups G admit S with \textbar S\textbar{} = Δ and diam H(G, S) ≤ D. (b)
If G = Q × \(Z_{m}\) with Q perfect (for example Q = SL(2, q) or PSL(2,
q), q ≥ 4) and diam H(G, S) ≤ D, then m ≤ A(Δ, D) = Θ(D\(^{Δ−1}\)) for
fixed Δ.

\emph{Proof.} (a) \textbar G\textbar{} ≤ M\(^{b}\)(Δ, D)/2 ({[}H{]} Thm.
4.1). (b) {[}H{]} Cor. 4.4(ii). ∎

So an infinite uniform family of the records' shape has Δ fixed and D →
∞ (and then q → ∞, since m grows at most polynomially in D while
\textbar G\textbar{} must grow exponentially), or D fixed and Δ → ∞. The
LPS family below can be read in both regimes (Corollary 3.4).

\begin{center}\rule{0.5\linewidth}{0.5pt}\end{center}

\hypertarget{the-lps-family-has-the-records-shape}{%
\subsection{2. The LPS family has the records'
shape}\label{the-lps-family-has-the-records-shape}}

\hypertarget{definitions}{%
\subsubsection{2.1 Definitions}\label{definitions}}

Let H(Z) = \{a₀ + a₁i + a₂j + a₃k : \(a_{t}\) ∈ Z\} be the Lipschitz
quaternions (i² = j² = k² = ijk = −1), with conjugate ᾱ = a₀ − a₁i − a₂j
− a₃k and norm N(α) = αᾱ = a₀² + a₁² + a₂² + a₃². N is multiplicative.

\textbf{Definition 2.1 (the LPS set).} For an odd prime p let \(S_{p}\)
be the set of α ∈ H(Z) with N(α) = p and

\begin{itemize}
\tightlist
\item
  if p ≡ 1 (mod 4): a₀ \textgreater{} 0 odd, and a₁, a₂, a₃ even;
\item
  if p ≡ 3 (mod 4): a₀ even, a₁, a₂, a₃ odd, and either a₀
  \textgreater{} 0, or a₀ = 0 and a₁ \textgreater{} 0.
\end{itemize}

By Jacobi's four-square theorem \textbar{}\(S_{p}\)\textbar{} = p + 1
(\emph{cited}: LPS, Combinatorica 8 (1988), §2, and G. Davidoff, P.
Sarnak, A. Valette, \emph{Elementary Number Theory, Group Theory and
Ramanujan Graphs}, LMS Student Texts 55 (2003), Ch. 2 {[}secondary{]}).
If α ∈ \(S_{p}\) then ᾱ or −ᾱ is in \(S_{p}\).

\textbf{Definition 2.2 (the splitting).} Let q be an odd prime, q ≠
p.~Choose x, y ∈ \(F_{q}\) with x² + y² + 1 = 0 (they exist: the sets
\{x²\} and \{−1 − y²\} both have (q + 1)/2 elements, so they meet). Put

\[\psi(a_0 + a_1 i + a_2 j + a_3 k) = \begin{pmatrix} a_0 + a_1 x + a_3 y & -a_1 y + a_2 + a_3 x \\ -a_1 y - a_2 + a_3 x & a_0 - a_1 x - a_3 y\end{pmatrix} \bmod q .\]

\textbf{Lemma 2.3.} ψ: H(Z) → \(M_{2}\)(\(F_{q}\)) is a ring
homomorphism with kernel qH(Z), and it is onto. det ψ(α) ≡ N(α), and
ψ(ᾱ) = adj ψ(α), so ψ(α)ψ(ᾱ) = N(α)·I.

\emph{Proof.} ψ is additive and ψ(1) = I. Put \(M_{i}\) = ψ(i) =
{[}{[}x, −y{]}, {[}−y, −x{]}{]}, \(M_{j}\) = ψ(j) = {[}{[}0, 1{]},
{[}−1, 0{]}{]}, \(M_{k}\) = ψ(k) = {[}{[}y, x{]}, {[}x, −y{]}{]}. Direct
multiplication gives \(M_{i}\)² = (x² + y²)·I = −I, \(M_{j}\)² = −I,
\(M_{i}\)·\(M_{j}\) = \(M_{k}\) and \(M_{j}\)·\(M_{i}\) = {[}{[}−y,
−x{]}, {[}−x, y{]}{]} = −\(M_{k}\). Since H(Z) is the ring generated
over Z by i, j subject to i² = j² = −1 and ji = −ij (with k = ij), the
assignment i ↦ \(M_{i}\), j ↦ \(M_{j}\) extends to a ring homomorphism,
and it sends k to \(M_{i}\)·\(M_{j}\) = \(M_{k}\); this is the displayed
formula. Its kernel contains qH(Z). If ψ(α) = 0, the diagonal gives a₀ +
a₁x + a₃y = a₀ − a₁x − a₃y = 0, so a₀ = 0 and a₁x + a₃y = 0; the
off-diagonal gives a₂ = 0 and a₃x − a₁y = 0. Then a₁(x² + y²) = x(a₁x +
a₃y) − y(a₃x − a₁y) = 0, so a₁ = 0 (x² + y² = −1), and then a₃x = a₃y =
0 with (x, y) ≠ (0, 0), so a₃ = 0, all mod q. So the kernel is qH(Z),
and the induced map H(Z)/qH(Z) → \(M_{2}\)(\(F_{q}\)) is injective
between sets of size q⁴, hence onto. The determinant of the displayed
matrix is (a₀ + a₁x + a₃y)(a₀ − a₁x − a₃y) − (a₂ + a₃x − a₁y)(−a₂ + a₃x
− a₁y) = a₀² − (a₁x + a₃y)² + a₂² − (a₃x − a₁y)² = a₀² + a₂² − (a₁² +
a₃²)(x² + y²) = N(α). Finally adj ψ(α) has diagonal (a₀ − a₁x − a₃y, a₀
+ a₁x + a₃y) and off-diagonal entries −(−a₁y + a₂ + a₃x), −(−a₁y − a₂ +
a₃x), which is ψ(a₀ − a₁i − a₂j − a₃k). ∎

\textbf{Lemma 2.4.} Let q \textgreater{} p.~(a) If α, β ∈ H(Z) have N(α)
= N(β) = p and ψ(α) = cψ(β) for some c ∈ \(F_{q}\)\^{}×, then α = ±β.
(b) The p + 1 points {[}ψ(α){]}, α ∈ \(S_{p}\), of PGL(2, q) are
distinct, and the set they form is closed under inversion.

\emph{Proof.} (a) γ = αβ̄ has N(γ) = p², and ψ(γ) = ψ(α)ψ(β̄) = cψ(β)ψ(β̄)
= cp·I. By Lemma 2.3, γ − c̃p ∈ qH(Z) for an integer c̃ ≡ c, so the
coordinates γ₁, γ₂, γ₃ are divisible by q. They are at most √N(γ) = p
\textless{} q in absolute value, so they vanish: γ is an integer with γ²
= p², γ = ±p.~Then α = γβ̄⁻¹ = ±pβ/N(β) = ±β. (b) Distinct elements of
\(S_{p}\) are not negatives of each other (the sign conditions of
Definition 2.1 fix the sign), so (a) gives distinctness. By Lemma 2.3,
{[}ψ(α){]}⁻¹ = {[}adj ψ(α){]} = {[}ψ(ᾱ){]} = {[}ψ(±ᾱ){]}, and ±ᾱ ∈
\(S_{p}\). ∎

\textbf{Definition 2.5'.} \(Γ_{p,q}\) is the connected component of (I,
A) in H(PGL(2, q), ψ(\(S_{p}\))).

\textbf{Lemma 2.5 (structure of \(Γ_{p,q}\)).} Let q \textgreater{} p
and let c ∈ ψ(\(S_{p}\)).

\begin{enumerate}
\def\labelenumi{(\alph{enumi})}
\tightlist
\item
  All of ψ(\(S_{p}\)) lies in one coset of PSL(2, q): in PSL(2, q) if p
  is a square mod q, in the other coset if not.
\item
  The A-vertices of \(Γ_{p,q}\) form H₀ × \{A\} for a subgroup H₀ ≤
  PSL(2, q), and its B-vertices cH₀ × \{B\}. The right translations
  \(ρ_{h}\), h ∈ H₀, map \(Γ_{p,q}\) to itself and act regularly on each
  side. When H₀ = PSL(2, q), \(Γ_{p,q}\) has n = q(q² − 1) vertices.
\item
  diam \(Γ_{p,q}\) = ecc(I, A) (computed in \(Γ_{p,q}\)).
\item
  \(Γ_{p,q}\) ≅ H(PSL(2, q), c⁻¹ψ(\(S_{p}\))) when it has q(q² − 1)
  vertices.
\end{enumerate}

\emph{Proof.} (a) det ψ(α) = p for all α (Lemma 2.3), and PSL(2, q) is
the kernel of the homomorphism PGL(2, q) →
\(F_{q}\)\textsuperscript{×/(\(F_{q}\)}×)² induced by the determinant.
(b) A letter and its inverse both have determinant class {[}p{]}, so an
A-word of length L has class {[}p{]}\(^{L}\): it lies in PSL(2, q) for L
even and in cPSL(2, q) for L odd. The set H₀ of the values of even
A-words (the A-side of \(Γ_{p,q}\)) is closed under products and
contains I, so it is a subgroup of the finite group PSL(2, q). Right
translation \(ρ_{h}\) by h ∈ H₀ is an automorphism ({[}H{]} Lemma 1.2)
mapping (I, A) to (h, A) ∈ \(Γ_{p,q}\), so it maps \(Γ_{p,q}\) to
itself; \{\(ρ_{h}\) : h ∈ H₀\} acts regularly on the A-side H₀ × \{A\}
and on the B-side cH₀ × \{B\}. When Theorem 2.7 applies, H₀ = PSL(2, q).
(c) By (b) and {[}H{]} Lemma 2.1 (with the group \{\(ρ_{h}\) : h ∈
H₀\}), diam \(Γ_{p,q}\) = max(ecc(I, A), ecc(c, B)). Since ψ(\(S_{p}\))
is closed under inversion (Lemma 2.4), the side swap σ(g, A) = (g, B),
σ(g, B) = (g, A) is an automorphism of H(PGL(2, q), ψ(\(S_{p}\)))
({[}H{]} Lemma 3.1(iv)); and \(ρ_{c}\) is one too. \(ρ_{c}\)∘σ maps (I,
A) to (c, B) and maps \(Γ_{p,q}\) onto the component of (c, B), which is
\(Γ_{p,q}\) since (c, B) \textasciitilde{} (I, A). Hence ecc(c, B) =
ecc(I, A). (d) The map (g, A) ↦ (g, A), (h, B) ↦ (c⁻¹h, B) is the
isomorphism H(G, S) → H(G, c⁻¹S) of {[}H{]} Lemma 3.1(i); restricted to
\(Γ_{p,q}\) it lands in PSL(2, q) × \{A, B\}, and both graphs are
Δ-regular with q(q² − 1) vertices. ∎

\hypertarget{the-order-3-symmetry}{%
\subsubsection{2.2 The order-3 symmetry}\label{the-order-3-symmetry}}

\textbf{Theorem 2.6.} Let p, q be odd primes, q \textgreater{} p, let ζ
= 1 + i + j + k and h = ψ(ζ).

\begin{enumerate}
\def\labelenumi{(\alph{enumi})}
\tightlist
\item
  h ∈ GL(2, q) has determinant 4, so {[}h{]} ∈ PSL(2, q), and {[}h{]}
  has order 3 in PGL(2, q).
\item
  For every α ∈ H(Z), hψ(α)h⁻¹ = ψ(σ(α)), where σ(a₀ + a₁i + a₂j + a₃k)
  = a₀ + a₃i + a₁j + a₂k.
\item
  Conjugation by h permutes the p + 1 points of ψ(\(S_{p}\)). Its fixed
  points are the {[}ψ(α){]} with a₁ = a₂ = a₃, and its cycle type is
\end{enumerate}

\begin{longtable}[]{@{}
  >{\raggedright\arraybackslash}p{(\columnwidth - 6\tabcolsep) * \real{0.2500}}
  >{\raggedright\arraybackslash}p{(\columnwidth - 6\tabcolsep) * \real{0.2500}}
  >{\raggedright\arraybackslash}p{(\columnwidth - 6\tabcolsep) * \real{0.2500}}
  >{\raggedright\arraybackslash}p{(\columnwidth - 6\tabcolsep) * \real{0.2500}}@{}}
\toprule\noalign{}
\begin{minipage}[b]{\linewidth}\raggedright
p
\end{minipage} & \begin{minipage}[b]{\linewidth}\raggedright
cycle type on ψ(\(S_{p}\))
\end{minipage} & \begin{minipage}[b]{\linewidth}\raggedright
Δ = p + 1
\end{minipage} & \begin{minipage}[b]{\linewidth}\raggedright
records with this cycle type ({[}H{]} §5.1)
\end{minipage} \\
\midrule\noalign{}
\endhead
\bottomrule\noalign{}
\endlastfoot
p = 3 & 1 + 3 & 4 & none (Δ = 4 records are metacyclic) \\
p ≡ 1 (mod 3) & 1 + 1 + 3\(^{(p−1)/3}\) & 8 (p = 7), 14 (p = 13), 20,
32, 38, \ldots{} & (8,8), (8,9), (8,10); (14,9), (14,10) \\
p ≡ 2 (mod 3) & 3\(^{(p+1)/3}\) & 6 (p = 5), 12 (p = 11), 18, 24, 30,
\ldots{} & (6,9), (6,10); (12,9), (12,10) \\
\end{longtable}

So the LPS Haar graph is an orbit-closed Haar graph in the sense of
{[}H{]} Construction 3.2, with α = conjugation by an element of order 3
of PSL(2, q) and m = 1, and with the same cycle type as the records at Δ
= 6, 8, 12 and 14.

\emph{Proof.} (a) N(ζ) = 4, so det h = 4 = 2², a nonzero square (q odd).
Write v = i + j + k, so v² = −3 and ζ = 1 + v; then ζ² = 1 + 2v + v² =
−2 + 2v and ζ³ = (−2 + 2v)(1 + v) = −2 − 2v + 2v + 2v² = −8. So h³ = −8I
is a scalar. And h is not a scalar: ψ(ζ) − c̃I = ψ(ζ − c̃) would vanish
only if ζ − c̃ ∈ qH(Z) (Lemma 2.3), but the i-coordinate of ζ − c̃ is 1.
So {[}h{]} has order 3.

\begin{enumerate}
\def\labelenumi{(\alph{enumi})}
\setcounter{enumi}{1}
\item
  We compute ζiζ̄. First ζi = i + i² + ji + ki = −1 + i + j − k. Then (−1
  + i + j − k)(1 − i − j − k) = (−1 + i + j + k) + (i + 1 + j − k) + (j
  + k + 1 − i) + (−k + j − i − 1) = 4j, using ij = k, ik = −j, ji = −k,
  jk = i, ki = j, kj = −i. So ζiζ̄ = 4j, that is, ζiζ⁻¹ = j. Applying the
  ring automorphism of H(Z) that permutes i → j → k → i (it preserves
  the relations and fixes ζ), we get ζjζ⁻¹ = k and ζkζ⁻¹ = i. Hence
  ζαζ⁻¹ = σ(α), and ζαζ̄ = 4σ(α). Applying ψ (Lemma 2.3): hψ(α)·adj(h) =
  4ψ(σ(α)), and adj(h) = det(h)·h⁻¹ = 4h⁻¹.
\item
  σ preserves the norm, the coordinate a₀ and the parity pattern of (a₁,
  a₂, a₃). If p ≡ 1 (mod 4), σ(\(S_{p}\)) = \(S_{p}\). If p ≡ 3 (mod 4),
  σ(α) satisfies Definition 2.1 except possibly the sign rule when a₀ =
  0, so σ(α) or −σ(α) lies in \(S_{p}\), and ψ(−β) = −ψ(β) is the same
  point of PGL(2, q). So conjugation by h permutes ψ(\(S_{p}\)). By
  Lemma 2.4(a), {[}ψ(σα){]} = {[}ψ(α){]} iff σ(α) = ±α. The case σ(α) =
  −α forces a₀ = 0 and (a₃, a₁, a₂) = (−a₁, −a₂, −a₃), so a₁ = −a₃ = a₂
  = −a₁, so α = 0, impossible. So the fixed points are the α ∈ \(S_{p}\)
  with a₁ = a₂ = a₃ =: a, i.e.~with a₀² + 3a² = p.
\end{enumerate}

\begin{itemize}
\tightlist
\item
  If p ≡ 2 (mod 3), a₀² + 3a² ≡ a₀² ∈ \{0, 1\} (mod 3), so there is no
  fixed point.
\item
  If p ≡ 1 (mod 3), the equation a₀² + 3a² = p has exactly four integer
  solutions (±a₀, ±a), with a₀, a ≠ 0 (Fermat; for a prime p ≡ 1 mod 3
  the representation p = x² + 3y² with x, y \textgreater{} 0 is unique;
  \emph{cited}: D. A. Cox, \emph{Primes of the form x² + ny²}, Wiley
  1989, §1 {[}secondary{]}). If p ≡ 1 (mod 4) then a₀ is odd and a even
  (odd + odd ≡ 0 and even + 3·odd ≡ 3 mod 4), so exactly the two
  solutions with a₀ \textgreater{} 0, namely a₀ + a(i + j + k) with a =
  ±\textbar a\textbar, are in \(S_{p}\). If p ≡ 3 (mod 4) then a₀ is
  even and a odd, a₀ ≠ 0 (p ≠ 3a²), and again exactly the two with a₀
  \textgreater{} 0 are in \(S_{p}\). So there are two fixed points.
\item
  If p = 3: a₀ = 0, a = ±1, and of ±(i + j + k) only i + j + k is in
  \(S_{3}\). One fixed point.
\end{itemize}

Since {[}h{]} has order 3, every other point lies in a 3-cycle, which
gives the cycle types. ∎

\textbf{Theorem 2.7 (\emph{cited}: Lubotzky--Phillips--Sarnak,
Combinatorica 8 (1988) 261--277 {[}secondary for the constants{]};
Davidoff--Sarnak--Valette (2003), Ch. 4 {[}secondary{]}).} Let p ≡ 1
(mod 4) and let q \textgreater{} p be an odd prime. Then \(Γ_{p,q}\) has
q(q² − 1) vertices (that is, the LPS Cayley graph X\(^{p,q}\) is
connected), every eigenvalue of its adjacency matrix other than ±(p + 1)
has absolute value at most 2√p (Ramanujan), and diam \(Γ_{p,q}\) ≤ 2
log\_p n + O(1), with an explicit constant given by LPS.

\emph{Remark.} \(Γ_{p,q}\) is X\(^{p,q}\) itself when p is not a square
mod q (X\(^{p,q}\) is then a bipartite Cayley graph of PGL(2, q), and a
bipartite Cayley graph whose connection set lies in the non-trivial
coset of an index-2 subgroup is a Haar graph of that subgroup, Lemma
2.5(d)); it is the bipartite double cover of X\(^{p,q}\) when p is a
square mod q. \emph{Computer-checked}: \(Γ_{p,q}\) has q(q² − 1)
vertices for p ∈ \{3, 5, 7, 11, 13\} and every prime p \textless{} q ≤
160 (logs \texttt{haar\_families/scan/lps\_p*.txt}, column
``vertices''), including p ≡ 3 (mod 4), which Theorem 2.7 as cited does
not cover.

\hypertarget{the-family-written-out-at-ux3b4-6}{%
\subsubsection{2.3 The family written out at Δ =
6}\label{the-family-written-out-at-ux3b4-6}}

\textbf{Example 2.8 (p = 5).} \(S_{5}\) = \{1 ± 2i, 1 ± 2j, 1 ± 2k\}.
With x² + y² = −1 in \(F_{q}\), the six matrices are

\[\psi(1 \pm 2i) = \begin{pmatrix} 1 \pm 2x & \mp 2y \\ \mp 2y & 1 \mp 2x\end{pmatrix},\quad
\psi(1 \pm 2j) = \begin{pmatrix} 1 & \pm 2 \\ \mp 2 & 1\end{pmatrix},\quad
\psi(1 \pm 2k) = \begin{pmatrix} 1 \pm 2y & \pm 2x \\ \pm 2x & 1 \mp 2y\end{pmatrix},\]

and conjugation by h = ψ(1 + i + j + k) = {[}{[}1 + x + y, 1 + x − y{]},
{[}−1 + x − y, 1 − x − y{]}{]} cycles 1 + 2i → 1 + 2j → 1 + 2k → 1 + 2i
and the same with −2: cycle type 3 + 3. In the coordinates of §1.2,
\textbf{the point {[}a : b : c : d{]} of the first copy of PSL(2, q) is
joined to the six points}

\[[\,(1 \pm 2x)a \mp 2yc : (1 \pm 2x)b \mp 2yd : \mp 2ya + (1 \mp 2x)c : \mp 2yb + (1 \mp 2x)d\,],\]
\[[\,a \pm 2c : b \pm 2d : c \mp 2a : d \mp 2b\,],\qquad
[\,(1 \pm 2y)a \pm 2xc : (1 \pm 2y)b \pm 2xd : \pm 2xa + (1 \mp 2y)c : \pm 2xb + (1 \mp 2y)d\,]\]

of the second copy; and the point {[}a : b : c : d{]} of the second copy
is joined to the same six expressions, read as points of the first copy
(by Lemma 2.3, adj ψ(α) = ψ(ᾱ), and S̄\_5 = \(S_{5}\), so the six
adjugates are the same six matrices up to scalars). For q = 13, with (x,
y) = (0, 5) (the first solution of x² + y² + 1 ≡ 0 found by the
programs), the component of the first copy's identity is a 6-regular
bipartite graph on 2184 vertices of diameter 7 (computer-checked by
\texttt{lps\_haar}, \texttt{haar\_ecc} and the all-pairs BFS
\texttt{apsp.c}, §8).

\begin{center}\rule{0.5\linewidth}{0.5pt}\end{center}

\hypertarget{the-43-obstruction}{%
\subsection{3. The 4/3 obstruction}\label{the-43-obstruction}}

\textbf{Lemma 3.1 (walk values).} Let p, q be odd primes, q
\textgreater{} p.~A walk of even length L in H(PGL(2, q), ψ(\(S_{p}\)))
from (I, A) ends at a vertex ({[}ψ(β){]}, A) with β ∈ H(Z), N(β) =
p\(^{L}\) and β ≡ 1 (mod 2H(Z)); in particular the real coordinate b₀ of
β is odd.

\emph{Proof.} The walk with letters ψ(\(α_{1}\)), \ldots, ψ(\(α_{L}\))
ends at \(s_{L}\)⁻¹s\_\{L−1\}⋯\(s_{2}\)⁻¹s\_1 with \(s_{t}\) =
ψ(\(α_{t}\)), and {[}ψ(α)⁻¹{]} = {[}ψ(ᾱ){]} (Lemma 2.3). So it ends at
{[}ψ(β){]} with β = ᾱ\_L \(α_{L−1}\) ⋯ ᾱ\_2 \(α_{1}\), and N(β) =
p\(^{L}\). Reduce mod 2: in H(Z)/2H(Z) we have −1 = 1, so i, j, k
commute there. If p ≡ 1 (mod 4), every \(α_{t}\) and ᾱ\_t is ≡ 1 (mod
2H(Z)), so β ≡ 1. If p ≡ 3 (mod 4), every \(α_{t}\) and ᾱ\_t is ≡ v := i
+ j + k, and v² = −3 ≡ 1, so β ≡ v\(^{L}\) = 1 since L is even. ∎

\textbf{Theorem 3.2 (distance to ψ(i)).} Let p, q be odd primes with q
\textgreater{} p, and let L₀ be the least even integer L ≥ 0 with
2p\(^{L/2}\) \textgreater{} q². In \(Γ_{p,q}\):

\begin{enumerate}
\def\labelenumi{(\alph{enumi})}
\tightlist
\item
  d((I, A), ({[}ψ(i){]}, A)) ≥ L₀, and the same holds for ψ(j) and ψ(k);
\item
  if p ≡ 1 (mod 4), equality holds.
\end{enumerate}

Note that det ψ(i) = N(i) = 1, so {[}ψ(i){]} ∈ PSL(2, q) for every q.

\emph{Proof.} (a) A path from (I, A) to ({[}ψ(i){]}, A) has even length
L (the graph is bipartite). By Lemma 3.1 there is β with N(β) =
p\(^{L}\), b₀ odd and {[}ψ(β){]} = {[}ψ(i){]}, i.e.~ψ(β) = cψ(i) with c
∈ \(F_{q}\)\^{}×. By Lemma 2.3, β − c̃i ∈ qH(Z), so q \textbar{} b₀, b₂,
b₃. Put P = p\(^{L/2}\), an integer. Since b₀ is odd, b₀ ≠ 0, so b₀² ≥
q², and then b₁² = P² − b₀² − b₂² − b₃² ≤ P² − q², so
\textbar b₁\textbar{} \textless{} P. Also q² divides b₀² + b₂² + b₃² =
P² − b₁² = (P − b₁)(P + b₁). The prime q does not divide P (q ≠ p), so q
cannot divide both factors (it would divide their sum 2P). Hence q²
divides P − εb₁ for some ε ∈ \{1, −1\}. Since \textbar b₁\textbar{}
\textless{} P, 0 \textless{} P − εb₁ \textless{} 2P, so q² ≤ P − εb₁
\textless{} 2P. That is 2p\(^{L/2}\) \textgreater{} q², and L ≥ L₀. For
ψ(j) and ψ(k): conjugation by h fixes (I, A), is an automorphism of
H(PGL(2, q), ψ(\(S_{p}\))) (Theorem 2.6 and {[}H{]} Prop. 3.3(a)), and
maps {[}ψ(i){]} to {[}ψ(j){]} to {[}ψ(k){]} (Theorem 2.6(b)), so the
three distances are equal.

\begin{enumerate}
\def\labelenumi{(\alph{enumi})}
\setcounter{enumi}{1}
\tightlist
\item
  Let L = L₀, P = p\(^{L/2}\), and put n₁ = 2P − q² \textgreater{} 0.
  Here P ≡ 1 (mod 4) (p ≡ 1 mod 4), so 2P ≡ 2 (mod 8), and q² ≡ 1 (mod
  8), so n₁ ≡ 1 (mod 8). By the Gauss--Legendre three-square theorem
  (\emph{cited}; n is a sum of three squares unless n = 4\(^{a}\)(8b +
  7)), n₁ = u² + v² + w² with integers u, v, w; since n₁ ≡ 1 (mod 4),
  exactly one of them is odd, and we name it u. Put
\end{enumerate}

\[\beta = qu + (P - q^2)\,i + qv\,j + qw\,k .\]

Then N(β) = q²(u² + v² + w²) + (P − q²)² = q²(2P − q²) + P² − 2Pq² + q⁴
= P² = p\(^{L}\); b₀ = qu is odd, b₁ = P − q² is even, and b₂ = qv, b₃ =
qw are even; so β ≡ 1 (mod 2H(Z)). And β ≡ P·i (mod qH(Z)), so
{[}ψ(β){]} = {[}Pψ(i){]} = {[}ψ(i){]}. By the LPS factorization theorem
(\emph{cited}: LPS 1988 §3; Davidoff--Sarnak--Valette (2003), Ch. 2, the
unique factorization of quaternions β ≡ 1 mod 2 of norm p\(^{L}\)
{[}secondary{]}), β = ±p\(^{r}\) \(α_{1}\)α\_2⋯\(α_{L−2r}\) with
\(α_{t}\) ∈ \(S_{p}\) and r ≥ 0. Since {[}ψ(±p\(^{r}\) γ){]} =
{[}ψ(γ){]}, we get {[}ψ(i){]} = {[}ψ(\(α_{1}\)⋯\(α_{L'}\)){]} with L' =
L − 2r even. Choose the letters \(s_{1}\) = ψ(\(α_{L'}\)), \(s_{2}\) =
ψ(ᾱ\_\{L'−1\}) (so {[}\(s_{2}\)⁻¹{]} = {[}ψ(\(α_{L'−1}\)){]}), \(s_{3}\)
= ψ(\(α_{L'−2}\)), \(s_{4}\) = ψ(ᾱ\_\{L'−3\}), \ldots; they lie in
ψ(\(S_{p}\)) up to sign (Lemma 2.4(b)), and the A-word \(s_{L'}\)⁻¹ ⋯
\(s_{2}\)⁻¹s\_1 has value {[}ψ(\(α_{1}\)⋯\(α_{L'}\)){]}. So there is a
walk of length L' ≤ L₀ from (I, A) to ({[}ψ(i){]}, A). With (a), the
distance is L₀. ∎

\textbf{Corollary 3.3.} For odd primes q \textgreater{} p, with n = q(q²
− 1) ≥ \textbar V(\(Γ_{p,q}\))\textbar:

\[\operatorname{diam}\Gamma_{p,q} \;\ge\; L_0 \;\ge\; 4\log_p q - 2\log_p 2 \;>\; \tfrac43 \log_p n - 2\log_p 2 .\]

\emph{Proof.} 2p\(^{L₀/2}\) \textgreater{} q² gives L₀ \textgreater{} 2
log\_p(q²/2) = 4 log\_p q − 2 log\_p 2, and n \textless{} q³. ∎

\textbf{Corollary 3.4 (a vanishing fraction of the Moore bound).} Let p
be an odd prime, D ≥ 2, and q \textgreater{} p a prime with diam
\(Γ_{p,q}\) ≤ D. Then q² \textless{} 2p\(^{⌊D/2⌋}\), the number of
vertices satisfies

\[n \le q(q^2-1) < 2^{3/2}\, p^{\frac32\lfloor D/2\rfloor} \le 2^{3/2} p^{3D/4},\]

and n \textless{} √2·p\(^{1−D/4}\)·M\(^{b}\)(p + 1, D). In particular,
along the LPS family the fraction of the Moore bound tends to 0 both for
fixed p and D → ∞ and for fixed D ≥ 5 and p → ∞.

\emph{Proof.} L₀ ≤ diam ≤ D and L₀ is even, so L₀ ≤ 2⌊D/2⌋ and
2p\(^{⌊D/2⌋}\) ≥ 2p\(^{L₀/2}\) \textgreater{} q². Then n \textless{} q³
\textless{} (2p\(^{⌊D/2⌋}\))\(^{3/2}\). Finally M\(^{b}\)(p + 1, D) =
\(2Σ_{j<D}\) p\(^{j}\) \textgreater{} 2p\(^{D−1}\), so n/M\(^{b}\)
\textless{} 2\(^{3/2}\)p\(^{3D/4}\)/(2p\(^{D−1}\)) = √2·p\(^{1−D/4}\). ∎

For comparison: at D = 8 the bound is √2/p (28 \% for p = 5), and the
LPS graphs actually reach 1--6 \% (Appendix B); the lab's records at D =
8 are at 8--22 \%, and the D = 8 table values with Δ ≥ 10 (Gómez's
compounds of generalized hexagons) at about 10 \% ({[}H{]} §5, §6). In
the regime D fixed, Δ → ∞, Delorme's constructions from generalized
polygons have lim inf n·Δ\(^{1−D}\) \textgreater{} 0 for every D (C.
Delorme, J. Graph Theory 9 (1985) 325--334 {[}abstract{]}), so
arithmetic Haar families lose to them by a factor p\(^{D/4−1}\).

\textbf{Computer-checked 3.5.} For p ∈ \{3, 5, 7, 11, 13, 17, 19, 29\}
and every prime q with p \textless{} q ≤ 160, the distance from (I, A)
to ({[}ψ(i){]}, A), ({[}ψ(j){]}, A) and ({[}ψ(k){]}, A), computed by
breadth-first search (\texttt{lps\_haar.c}), equals L₀: no exception in
251 cases (log \texttt{haar\_families/recheck\_thmB.log}). This includes
p ≡ 3 (mod 4), where equality is not proved here. The eccentricity of
(I, A) exceeds L₀ by 0, 1, 2 or 3 (in 76, 151, 22 and 2 of the 251
cases). So in this range the special vertex ψ(i) accounts for the
diameter up to at most 3, and at odd D it is usually what limits the
family (Appendix B). The construction in the proof of Theorem 3.2(b) (u,
v, w, and β with N(β) = p\(^{L₀}\), β ≡ 1 mod 2 and {[}ψ(β){]} =
{[}ψ(i){]}) was also checked directly for p ∈ \{5, 13, 17, 29, 37\} and
all primes p \textless{} q \textless{} 1000 (802 cases,
\texttt{haar\_families/check\_thmB\_construction.py}). An independent
re-implementation by squad card c8 (its own quaternion arithmetic and
BFS; \texttt{research/squads/thmfam-r1/c8/}) confirmed Theorem 2.6(c)
for all 137 pairs with 7 ≤ q ≤ 61 and Theorem 3.2 in 41 cases with q ≡ 1
(mod 4), and its all-pairs BFS (\texttt{apsp.c}) agreed with
\texttt{haar\_ecc} on six LPS graphs. (The squad card also found a bug
in the first version of \texttt{lps\_haar.c}, which mixed the distances
of the two sides; the logs above are from the corrected version.)

\textbf{Proposition 3.6 (the same obstruction in other definite
quaternion orders).} Let 𝒪 be a Z-order in a definite quaternion algebra
over Q, with norm form N, and suppose 𝒪 = Zγ ⊕ 𝒪′ with N(aγ + β′) = a² +
N(β′) for a ∈ Z, β′ ∈ 𝒪′ (an orthogonal unit γ, N(γ) = 1). Let ψ: 𝒪 →
\(M_{2}\)(\(F_{q}\)) be a ring homomorphism with kernel q𝒪 (q an odd
prime, q ∤ p), let S ⊆ 𝒪 be a set of elements of norm p closed under
conjugation up to sign, and suppose that no product of an even number L
≥ 2 of elements of S lies in Zγ. Then in H(PGL(2, q), ψ(S)), every walk
of even length L from (I, A) to ({[}ψ(γ){]}, A) has 2p\(^{L/2}\)
\textgreater{} q².

\emph{Proof.} As in Lemma 3.1 the walk ends at {[}ψ(β){]} with β a
product of L elements of S (conjugates included), so N(β) = p\(^{L}\);
{[}ψ(β){]} = {[}ψ(γ){]} gives β = bγ + qβ″ with b ∈ Z, β″ ∈ 𝒪′ (the
kernel of ψ is q𝒪 and 𝒪 = Zγ ⊕ 𝒪′), and β″ ≠ 0 by hypothesis. Then P² =
b² + q²N(β″) with P = p\(^{L/2}\) and N(β″) ≥ 1, and the second half of
the proof of Theorem 3.2(a) applies verbatim (q ∤ P, q² \textbar{} P −
εb, 0 \textless{} P − εb \textless{} 2P). ∎

For the Hamilton order, γ = i, 𝒪′ = Z + Zj + Zk, and the hypothesis is
Lemma 3.1 (every even product is ≡ 1 mod 2, while elements of Zi have
even real part). \emph{Computer-checked} (squad card c9 of round 2,
\texttt{research/squads/thmfam-r2/c9/}): in the order Z⟨1, i, j, ij⟩
with i² = −1, j² = −3 (norm a₀² + a₁² + 3a₂² + 3a₃²), with S a set of p
+ 1 classes of elements of norm p modulo the units ±1, ±i, for p ∈ \{5,
7, 11\} and q ≤ 97, the distance from (I, A) to ({[}ψ(i){]}, A) equals
the least even L with 2p\(^{L/2}\) \textgreater{} q² in 65 of 66 cases
(in the remaining case, p = 5, q = 61, it is 12 instead of 10), and ψ(i)
is at maximal distance for p = 5, q ∈ \{61, 83, 97\}. So the obstruction
is a property of definite quaternion arithmetic with a unit orthogonal
to the rest, not of Hamilton's order. (The card did not check the
hypothesis ``no even product in Zi'' for its set; the lower bound it
measured holds in all 66 cases.)

\textbf{Remark 3.7 (prior art).} The lower bound (4/3) log\_p n + O(1)
for the diameter of the bipartite LPS graphs is due to N. T. Sardari,
\emph{Diameter of Ramanujan graphs and random Cayley graphs},
Combinatorica 39 (2019) 427--446, arXiv:1511.09340 {[}abstract{]}; his
abstract also states that almost all pairs of vertices of a Ramanujan
graph are at distance at most (1 + ε) \(log_{p}\) n, and reports
numerics in which LPS diameters grow like (4/3) log and those of random
Cayley graphs like log. We did not read his proof; the argument above is
short and may be his. The exact value in Theorem 3.2(b) follows from the
same arithmetic and the LPS factorization; we did not check whether it
is stated there. That the LPS sets are permuted by the Hurwitz unit (1 +
i + j + k)/2 is classical; we found no source that compares their cycle
types with orbit-searched connection sets, and claim nothing more than
the observation.

\begin{center}\rule{0.5\linewidth}{0.5pt}\end{center}

\hypertarget{cyclic-factors}{%
\subsection{4. Cyclic factors}\label{cyclic-factors}}

\textbf{Proposition 4.1.} Let K be a finite group, κ: \(S_{p}\) → K any
map, and S̃ = \{(ψ(α), κ(α)) : α ∈ \(S_{p}\)\} ⊆ PGL(2, q) × K. Then for
every k ∈ K, the distance in H(PGL(2, q) × K, S̃) from (I, \(e_{K}\), A)
to ({[}ψ(i){]}, k, A) is at least L₀. Hence if the component of (I,
\(e_{K}\), A) has diameter ≤ D and contains ({[}ψ(i){]}, k, A) for some
k, then q² \textless{} 2p\(^{⌊D/2⌋}\).

\emph{Proof.} The projection (g, k, X) ↦ (g, X) maps every edge \{(g, k,
A), (ψ(α)g, κ(α)k, B)\} to the edge \{(g, A), (ψ(α)g, B)\}, so it maps
walks to walks of the same length and does not decrease distances. Apply
Theorem 3.2(a). ∎

So a cyclic factor cannot help at the vertex ψ(i) with a fixed q; what
it changes is that m is free. The order is q(q² − 1)·m, the obstruction
bounds only q, and the counting bound ({[}H{]} Thm. 4.1) allows m up to
about p\(^{D−1}\)/q³. Whether the cyclic part can be filled efficiently
is a computational question:

\textbf{Computer-checked 4.2 (squad cards c6, c7).} Let r have order 3
mod m. Labels α ↦ z(α) ∈ \(Z_{m}\) that are equivariant for the order-3
symmetry (z(σα) = r·z(α); fixed letters get 0; one 3-cycle gets (1, r,
r²) without loss of generality, the others (t, rt, r²t)) make H(PGL(2,
q) × \(Z_{m}\), S̃) orbit-closed with α(x, z) = (hxh⁻¹, rz), exactly the
records' shape. Searching t for each (q, m), the squads found:

\begin{longtable}[]{@{}
  >{\raggedright\arraybackslash}p{(\columnwidth - 10\tabcolsep) * \real{0.1364}}
  >{\raggedright\arraybackslash}p{(\columnwidth - 10\tabcolsep) * \real{0.1364}}
  >{\raggedleft\arraybackslash}p{(\columnwidth - 10\tabcolsep) * \real{0.1818}}
  >{\raggedleft\arraybackslash}p{(\columnwidth - 10\tabcolsep) * \real{0.1818}}
  >{\raggedleft\arraybackslash}p{(\columnwidth - 10\tabcolsep) * \real{0.1818}}
  >{\raggedleft\arraybackslash}p{(\columnwidth - 10\tabcolsep) * \real{0.1818}}@{}}
\toprule\noalign{}
\begin{minipage}[b]{\linewidth}\raggedright
cell
\end{minipage} & \begin{minipage}[b]{\linewidth}\raggedright
best LPS × \(Z_{m}\) (q, m)
\end{minipage} & \begin{minipage}[b]{\linewidth}\raggedleft
vertices
\end{minipage} & \begin{minipage}[b]{\linewidth}\raggedleft
of table
\end{minipage} & \begin{minipage}[b]{\linewidth}\raggedleft
lab record
\end{minipage} & \begin{minipage}[b]{\linewidth}\raggedleft
LPS alone (Appendix B)
\end{minipage} \\
\midrule\noalign{}
\endhead
\bottomrule\noalign{}
\endlastfoot
(6,8) & (13, 13) & 28,392 & 95.3 \% & 31,680 & 12,144 \\
(6,9) & (13, 43) & 93,912 & 80.0 \% & 127,296 & 24,360 \\
(6,10) & (23, 37) & 449,328 & 99.4 \% & 524,160 & 148,824 \\
(8,8) & (11, 151) & 199,320 & 98.9 \% & 239,904 & 29,760 \\
(8,9) & (23, 73) & 886,512 & 81.2 \% & 1,341,600 & 300,696 \\
(14,8) & (29, 379) & 9,232,440 & 65.1 \% & none (the lab does not beat
this cell) & 1,294,920 \\
\end{longtable}

Every entry was re-checked by the director with \texttt{haar\_ecc}
(two-source BFS on the component, from the cards' generators:
\texttt{research/squads/thmfam-r1/c6/gen\_c6.py},
\texttt{research/squads/thmfam-r1/c7/c7lib.py}), and (6,10) and (8,8)
also with \texttt{verify\_cover.py} (independent Python arithmetic;
eccentricities 10, 10 and 8, 8). These graphs are below the lab's
records and below the table, and their labels were searched for each (q,
m): they are not a uniform family. They show that the cyclic factor
removes most of the 4/3 loss in practice: at (6,10) the cover is 3.0
times the LPS graph alone.

\begin{center}\rule{0.5\linewidth}{0.5pt}\end{center}

\hypertarget{other-uniform-families-squad-round-1}{%
\subsection{5. Other uniform families (squad round
1)}\label{other-uniform-families-squad-round-1}}

\textbf{Score.} For a family S(q) and the primes 23 ≤ q ≤ 113, the
\emph{excess} is diam − \(log_{Δ−1}\)\textbar G\textbar{} with
\textbar G\textbar{} = one side of the graph, averaged over q
(\texttt{haar\_families/family\_score.py}, which runs
\texttt{haar\_ecc}). Lower is better. A uniform family gets one S per q,
so it should be compared with \emph{one} random orbit-closed draw per q,
not with the records, which are the best of millions of evaluations for
one (q, m).

\textbf{Computer-checked 5.1 (mean excess, 22 primes 23 ≤ q ≤ 113, group
PSL(2, q) unless stated; from
\texttt{research/squads/thmfam-r1/scoreboard.jsonl} and the cards'
folders).} Baselines first (card c13: 8 seeds each, mean ± standard
deviation over seeds):

\begin{longtable}[]{@{}
  >{\raggedright\arraybackslash}p{(\columnwidth - 6\tabcolsep) * \real{0.2000}}
  >{\raggedleft\arraybackslash}p{(\columnwidth - 6\tabcolsep) * \real{0.2667}}
  >{\raggedleft\arraybackslash}p{(\columnwidth - 6\tabcolsep) * \real{0.2667}}
  >{\raggedleft\arraybackslash}p{(\columnwidth - 6\tabcolsep) * \real{0.2667}}@{}}
\toprule\noalign{}
\begin{minipage}[b]{\linewidth}\raggedright
one random draw per q
\end{minipage} & \begin{minipage}[b]{\linewidth}\raggedleft
Δ = 6
\end{minipage} & \begin{minipage}[b]{\linewidth}\raggedleft
Δ = 7
\end{minipage} & \begin{minipage}[b]{\linewidth}\raggedleft
Δ = 8
\end{minipage} \\
\midrule\noalign{}
\endhead
\bottomrule\noalign{}
\endlastfoot
random S, no symmetry, PSL & 3.20 ± 0.05 & 2.95 ± 0.05 & 2.95 ± 0.05 \\
random S, no symmetry, SL & 3.05 ± 0.02 & 2.89 ± 0.03 & \\
random orbit-closed S (k = 3), PSL & 3.31 ± 0.11 & 3.04 ± 0.09 & 3.28 ±
0.08 \\
random orbit-closed S (k = 3), SL & 3.34 ± 0.10 & 3.09 ± 0.08 & \\
\end{longtable}

So a single orbit-closed draw is slightly \emph{worse} than a single
unconstrained draw (by 0.1--0.3); the orbit search wins because it makes
the search space smaller and the misses come in clumps ({[}H{]} §3.5),
not because orbit-closed sets are better one by one. The families:

\begin{longtable}[]{@{}
  >{\raggedright\arraybackslash}p{(\columnwidth - 6\tabcolsep) * \real{0.2308}}
  >{\raggedright\arraybackslash}p{(\columnwidth - 6\tabcolsep) * \real{0.2308}}
  >{\raggedleft\arraybackslash}p{(\columnwidth - 6\tabcolsep) * \real{0.3077}}
  >{\raggedright\arraybackslash}p{(\columnwidth - 6\tabcolsep) * \real{0.2308}}@{}}
\toprule\noalign{}
\begin{minipage}[b]{\linewidth}\raggedright
family (one S per q, no search per q)
\end{minipage} & \begin{minipage}[b]{\linewidth}\raggedright
Δ
\end{minipage} & \begin{minipage}[b]{\linewidth}\raggedleft
mean excess
\end{minipage} & \begin{minipage}[b]{\linewidth}\raggedright
card
\end{minipage} \\
\midrule\noalign{}
\endhead
\bottomrule\noalign{}
\endlastfoot
LPS, p = 5 (component in PGL(2, q)) & 6 & 4.00 & director \\
integer lifts, O(X) ∪ O(Y), X, Y ∈ SL(2, Z), entries ≤ 3 (best of 206) &
6 & 3.36 & c1 \\
Dickson / trace / power formulas (best of 15) & 6 & 3.13 & c11 \\
integer lifts, \{I\} ∪ O(X) ∪ O(Y) (best of 1212; mean of the top 10:
3.35) & 7 & 3.10 & c2 \\
primitive-root formulas, e.g.~X = {[}{[}g, 1{]}, {[}0, g⁻¹{]}{]}, Y =
{[}{[}1/2, g{]}, {[}g, 2(1 + g²){]}{]} (best of 16) & 7 & 2.96 & c10 \\
integer lifts × \(Z_{13}\), equivariant labels (best of about 500
screened) & 7 & 3.12 & c4 \\
integer lifts, two fixed letters + two orbits (all 234 with entries ≤ 2)
& 8 & 3.84 & c3 \\
Dickson / trace / power formulas (best of 15) & 8 & 3.16 & c11 \\
conic model: Cayley transforms of integer matrices (best of 8) & 8 &
3.34 & c16 \\
integer lifts, one fixed letter + three orbits (all 858 with entries ≤
2) & 10 & 3.40 (random orbit-closed: 2.76) & c3 \\
single orbits of order 6 (best uniform x, y) & 12 & 2.73 (random, k = 3:
2.55--2.64) & c15 \\
LPS + one extra α-orbit (repair) & 9 & 3.24 on 7 primes (random:
2.72--2.86) & c18 \\
the records' own sets, lifted to SL(2, Z) and used at the other 21
primes & 6, 7 & 2.87--3.24 (at their own q: 1.99--2.13) & c5 \\
SL(2, q) × \(Z_{m}\), m ≈ q, integer letters with equivariant labels (10
primes 23--61) & 6 & 3.38 (same letters, m = 1: 3.60) & c17 \\
the lab's records (one (q, m) each, annealed) & 6--16 & about 2.0--2.3 &
{[}H{]} \\
\end{longtable}

Torus-orbit families with Δ growing like q (card c12: S = \{I\} ∪
T-conjugates of x for a torus T of PSL(2, q)) have diameter 5 to 7 for
31 ≤ q ≤ 61 and reach at most 1.8 \% of the Moore bound at their
diameter; no family of diameter ≤ 4 appeared. The excess of LPS grows
with log\_5\textbar G\textbar{} with slope 0.41 (card c14, 38 primes 23
≤ q ≤ 199; Corollary 3.3 predicts 1/3), that of random orbit-closed
draws with slope −0.04.

\textbf{Heuristic 5.2 (reading of 5.1).}

\begin{enumerate}
\def\labelenumi{\arabic{enumi}.}
\tightlist
\item
  Every formula family lands within about 0.3 of one random draw at the
  same Δ, on either side; the best of each kind was selected among
  8--1212 variants on the same primes, so the small wins are within the
  selection effect.
\item
  Arithmetic structure costs: LPS is about 0.8 worse than one random
  draw at q ≤ 113, and the gap grows like (1/3) log\_p n (Theorem 3.2;
  card c14 measured slope 0.41 for LPS and −0.04 for random draws).
\item
  The records do not transport: the set that is a record at its own q
  (excess about 2.0) is an ordinary set at every other q (card c5). A
  record is a property of the pair (q, S), not of a formula.
\item
  Larger symmetry (k = 5, 6) is worse than k = 3 at the same Δ, for
  formulas and for random draws (card c15).
\item
  The records are about 1 below a single draw. In the model of {[}H{]}
  §3.5 a single random S misses about μ = N e\(^{−W/N}\) elements and is
  valid with probability about e\(^{−μ}\); a search over T independent
  sets can afford μ* ≈ ln T, which lowers the needed W/N by about ln μ*,
  roughly one unit of \(log_{Δ−1}\) at the sizes of the records. A
  formula gives one set per q, so it cannot have this advantage unless
  it is better than random for structural reasons, as the Singer sets
  are at D = 3. We found no such structure at D ≥ 7.
\end{enumerate}

\begin{center}\rule{0.5\linewidth}{0.5pt}\end{center}

\hypertarget{why-no-record-quality-uniform-family-is-in-sight}{%
\subsection{6. Why no record-quality uniform family is in
sight}\label{why-no-record-quality-uniform-family-is-in-sight}}

\textbf{6.1 Fixed Δ, D → ∞.} By Proposition 1.1 a family here has q → ∞.
A uniform family of record quality would have diam ≤ \(log_{Δ−1}\) n +
O(log log n) along an explicit sequence of groups. No explicit family of
bounded-degree graphs with diam ≤ (1 + ε) \(log_{Δ−1}\) n is known
{[}secondary{]}; random Δ-regular graphs have diameter (1 + o(1))
\(log_{Δ−1}\) n (B. Bollobás, W. Fernandez de la Vega, Combinatorica 2
(1982) 125--134 {[}secondary{]}), and so, conjecturally, do random Haar
graphs of SL(2, q) (research/theory/THRESHOLD.md, Conjecture 6, on
master). The best explicit candidates, the Ramanujan graphs, have almost
all distances (1 + o(1)) \(log_{Δ−1}\) n (Sardari {[}abstract{]}; E.
Lubetzky, Y. Peres, \emph{Cutoff on all Ramanujan graphs}, GAFA 26
(2016) {[}secondary{]}), but for the LPS graphs a few vertices are at
distance (4/3) log\_p n (Theorem 3.2). That is the precise way in which
the most natural uniform family of the records' shape fails, and
Corollary 3.4 turns it into a loss of a factor p\(^{D/4−1}\) against the
Moore bound.

\textbf{6.2 Fixed D, Δ → ∞.} The LPS family with q ≈ p\(^{D/4}\) has
degree p + 1 and order about p\(^{3D/4}\) at diameter ≈ D (Corollary 3.4
is an upper bound; the lower side is Theorem 2.7, order
p\(^{(D−O(1))/2}\)). Delorme's polygon constructions reach a fixed
fraction of p\(^{D−1}\) (§3). We know of no Haar-graph family of the
records' groups that reaches lim inf n·Δ\(^{1−D}\) \textgreater{} 0 at D
= 8, 9 or 10; the random-S theorem of research/theory/THRESHOLD.md
reaches n ≥ Δ\(^{D−1}\)/((D − 1)!·(1 + η)·ln n) for every finite group,
which loses a logarithm.

\textbf{6.3 What would change this.} A formula that beats one random
draw by about 1 in excess, uniformly in q; or a proof that some explicit
cover (for example LPS × \(Z_{m}\) with labels given by a rule) has
diameter log\_p n + O(log log n). Neither is in sight. The squads' round
2 looked for label rules for LPS × \(Z_{m}\) (Appendix A).

\begin{center}\rule{0.5\linewidth}{0.5pt}\end{center}

\hypertarget{prior-art}{%
\subsection{7. Prior art}\label{prior-art}}

\begin{itemize}
\tightlist
\item
  \textbf{LPS graphs.} A. Lubotzky, R. Phillips, P. Sarnak,
  \emph{Ramanujan graphs}, Combinatorica 8 (1988) 261--277; G. A.
  Margulis, Problems Inform. Transmission 24 (1988) 39--46
  {[}secondary{]}; the book of Davidoff--Sarnak--Valette (2003)
  {[}secondary{]}. Definitions 2.1--2.2, Lemma 2.3, Theorem 2.7 and the
  factorization used in Theorem 3.2(b) are theirs.
\item
  \textbf{The 4/3 lower bound.} Sardari (2019) {[}abstract{]}, Remark
  3.7.
\item
  \textbf{Typical distances in Ramanujan graphs.} Lubetzky--Peres (2016)
  {[}secondary{]}; Sardari (2019) {[}abstract{]}.
\item
  \textbf{Haar graphs of SL(2, q) and PSL(2, q) in extremal problems.}
  {[}H{]} §6.1: Exoo--Jajcay (girth, a Haar graph of SL(2, 83));
  Bray--Parker--Rowley (the group PSL(2, q) × \(Z_{3}\)); the LPS graphs
  themselves are Haar graphs of PSL(2, q) in the bipartite case (Lemma
  2.5(d)). In the bipartite degree--diameter table at D = 8--10
  (Miller--Širáň, DS14v2, 2013, as reproduced by the Combinatorics Wiki)
  we found no LPS graph; the novelty check
  \texttt{research/lit/novelty-2026-10-25.md} (branch
  \texttt{claude/novelty-literature-check-4i5frp}, claim 5) found no
  Haar graphs of SL(2, q) or PSL(2, q) there.
\item
  \textbf{Asymptotic lower bounds for N\(^{b}\)(Δ, D).} Delorme (1985)
  {[}abstract via the novelty check{]}; random bipartite graphs (V.
  Klee, D. Larman, E. M. Wright, \emph{The diameter of almost all
  bipartite graphs}, Studia Sci. Math. Hungar. 15 (1980) 39--43
  {[}secondary{]}); research/theory/THRESHOLD.md for random Haar graphs.
\item
  \textbf{Orbit-closed connection sets.} {[}H{]} §6.1 (an old device;
  Singer sets and their multipliers).
\end{itemize}

What this note adds: Theorem 2.6 (the LPS sets have the records' orbit
shape, with the cycle type computed), the elementary proof and the exact
value in Theorem 3.2, Corollary 3.4 in the form of a fraction of the
Moore bound, the covers of §4 and the measurements of §5. We claim no
novelty for Lemmas 2.3--2.5 or for Theorem 2.7.

\begin{center}\rule{0.5\linewidth}{0.5pt}\end{center}

\hypertarget{how-to-re-check}{%
\subsection{8. How to re-check}\label{how-to-re-check}}

\begin{Shaded}
\begin{Highlighting}[]
\BuiltInTok{cd}\NormalTok{ research/theorems/haar\_families}
\FunctionTok{cc} \AttributeTok{{-}O3} \AttributeTok{{-}o}\NormalTok{ lps\_haar lps\_haar.c }\AttributeTok{{-}lm}      \CommentTok{\# LPS graph: generators ({-}v), eccentricities, d(I, psi(i)) etc.}
\FunctionTok{cc} \AttributeTok{{-}O3} \AttributeTok{{-}o}\NormalTok{ haar\_ecc haar\_ecc.c          }\CommentTok{\# any Haar graph of SL/PSL/PGL(2,q) x Z\_m from generator lines "a b c d z"}
\ExtensionTok{./lps\_haar}\NormalTok{ 5 13 }\AttributeTok{{-}v}                     \CommentTok{\# Example 2.8: 2184 vertices, diameter 7, d(I, psi(i)) = 6 = L0}
\ExtensionTok{python3}\NormalTok{ check\_thmB.py 3,5,7,11,13,17,19,29 160 }\OperatorTok{\textgreater{}}\NormalTok{ recheck\_thmB.log     }\CommentTok{\# Computer{-}checked 3.5}
\ExtensionTok{python3}\NormalTok{ check\_thmB\_construction.py                                    }\CommentTok{\# the beta of Theorem 3.2(b), q \textless{} 1000}
\ExtensionTok{./lps\_haar}\NormalTok{ 5 13 }\AttributeTok{{-}v} \KeywordTok{|} \FunctionTok{grep} \StringTok{\textquotesingle{}\^{}gen\textquotesingle{}} \KeywordTok{|} \FunctionTok{sed} \StringTok{\textquotesingle{}s/.*mat(//; s/) z=/ /; s/,/ /g\textquotesingle{}} \OperatorTok{\textgreater{}}\NormalTok{ g.txt}
\ExtensionTok{python3}\NormalTok{ component\_edges.py pgl 13 1 e.txt }\OperatorTok{\textless{}}\NormalTok{ g.txt                    }\CommentTok{\# independent edge list of the component}
\FunctionTok{cc} \AttributeTok{{-}O2} \AttributeTok{{-}o}\NormalTok{ apsp ../../lifts/coordinator\_check/apsp.c }\KeywordTok{\&\&} \ExtensionTok{./apsp}\NormalTok{ e.txt 6 40   }\CommentTok{\# all{-}pairs BFS: diameter 7}
\ExtensionTok{python3}\NormalTok{ table\_lps.py                   }\CommentTok{\# Appendix B from scan/lps\_p*.txt}
\ExtensionTok{python3}\NormalTok{ family\_score.py }\AttributeTok{{-}{-}gen} \StringTok{"./gen\_lps.sh 5 \{q\}"} \AttributeTok{{-}{-}group}\NormalTok{ pgl }\AttributeTok{{-}{-}qrange}\NormalTok{ 23 113   }\CommentTok{\# §5, LPS row}
\ExtensionTok{python3}\NormalTok{ ../../squads/thmfam{-}r1/c6/gen\_c6.py 23 37 19 0 }\KeywordTok{|} \ExtensionTok{./haar\_ecc}\NormalTok{ pgl 23 37  }\CommentTok{\# §4.2, (6,10)}
\ExtensionTok{python3}\NormalTok{ ../../squads/thmfam{-}r1/c6/gen\_c6.py 23 37 19 0 }\KeywordTok{|} \ExtensionTok{python3}\NormalTok{ verify\_cover.py 23 37}
\end{Highlighting}
\end{Shaded}

\begin{center}\rule{0.5\linewidth}{0.5pt}\end{center}

\hypertarget{appendix-a.-the-squad-rounds}{%
\subsection{Appendix A. The squad
rounds}\label{appendix-a.-the-squad-rounds}}

Two rounds of Claude Haiku cards (the lab's \texttt{haiku-squad} skill;
model haiku, effort high), directed by this session. The cards, the
shared rules and every card's code and notes are in
\texttt{research/squads/thmfam-r1/} and
\texttt{research/squads/thmfam-r2/}, one line per card in
\texttt{scoreboard.jsonl}; the rounds are logged in
\texttt{research/squads/LOG.md}. The squads searched constructions and
tested the theorems on small q; the proofs are the director's. Every
number quoted in the note from a card was re-run by the director with
the commands of §8, except where it says ``card'' only (the mean
excesses of §5, which are averages over many runs and are reproducible
from the cards' folders).

\textbf{Round 1 (18 cards): uniform families.} c1--c4 integer lifts (Δ =
6, 7, 8, 10, and with \(Z_{m}\)); c5 the records' sets at other q;
c6--c7 LPS × \(Z_{m}\) with searched labels; c8 verification of Theorems
2.6 and 3.2 (found a bug in the first \texttt{lps\_haar.c}); c9 the far
vertices of LPS graphs (the guess ``all far vertices come from pure
quaternions'' failed); c10--c11 q-dependent formulas; c12 torus orbits
with Δ \textasciitilde{} q; c13--c14 random baselines and the trend in
q; c15 symmetries of order 5 and 6; c16 the conic model; c17 cyclic
factors with m \textasciitilde{} q; c18 repairs of LPS. Outcome: every
formula family is at the level of one random draw (§5); LPS × \(Z_{m}\)
covers were the only direction that came close to the table (§4).

\textbf{Round 2 (11 cards): LPS × \(Z_{m}\) covers.} {[}{[}R2{]}{]}

\hypertarget{appendix-b.-the-lps-family-cell-by-cell}{%
\subsection{Appendix B. The LPS family cell by
cell}\label{appendix-b.-the-lps-family-cell-by-cell}}

\textbf{Computer-checked} (\texttt{haar\_families/appendix\_b.py} from
the scan logs \texttt{haar\_families/scan/lps\_p*.txt}, which list the
order and the diameter of \(Γ_{p,q}\) for every prime p \textless{} q up
to the scan limit, computed by \texttt{lps\_haar}). For each cell the
largest LPS graph of diameter ≤ D; ``(scan limit)'' means the largest
scanned q still had diameter ≤ D, so larger LPS graphs may exist there.
``Thm 3.2: q ≤'' is the largest prime q with q² \textless{}
2p\(^{⌊D/2⌋}\), the limit of Corollary 3.4.

\begin{longtable}[]{@{}
  >{\raggedright\arraybackslash}p{(\columnwidth - 14\tabcolsep) * \real{0.1034}}
  >{\raggedright\arraybackslash}p{(\columnwidth - 14\tabcolsep) * \real{0.1034}}
  >{\raggedright\arraybackslash}p{(\columnwidth - 14\tabcolsep) * \real{0.1034}}
  >{\raggedleft\arraybackslash}p{(\columnwidth - 14\tabcolsep) * \real{0.1379}}
  >{\raggedleft\arraybackslash}p{(\columnwidth - 14\tabcolsep) * \real{0.1379}}
  >{\raggedleft\arraybackslash}p{(\columnwidth - 14\tabcolsep) * \real{0.1379}}
  >{\raggedleft\arraybackslash}p{(\columnwidth - 14\tabcolsep) * \real{0.1379}}
  >{\raggedleft\arraybackslash}p{(\columnwidth - 14\tabcolsep) * \real{0.1379}}@{}}
\toprule\noalign{}
\begin{minipage}[b]{\linewidth}\raggedright
Δ (p)
\end{minipage} & \begin{minipage}[b]{\linewidth}\raggedright
D
\end{minipage} & \begin{minipage}[b]{\linewidth}\raggedright
largest LPS graph: q, vertices
\end{minipage} & \begin{minipage}[b]{\linewidth}\raggedleft
\% of Moore
\end{minipage} & \begin{minipage}[b]{\linewidth}\raggedleft
Thm 3.2: q ≤
\end{minipage} & \begin{minipage}[b]{\linewidth}\raggedleft
table
\end{minipage} & \begin{minipage}[b]{\linewidth}\raggedleft
ratio table/LPS
\end{minipage} & \begin{minipage}[b]{\linewidth}\raggedleft
lab record
\end{minipage} \\
\midrule\noalign{}
\endhead
\bottomrule\noalign{}
\endlastfoot
4 (3) & 7 & 7, 336 & 15.37 & 7 & 840 & 2.5 & 864 \\
4 (3) & 8 & 7, 336 & 5.12 & 11 & 2,184 & 6.5 & -- \\
4 (3) & 9 & 11, 1,320 & 6.71 & 11 & 4,970 & 3.8 & -- \\
4 (3) & 10 & 13, 2,184 & 3.70 & 19 & 11,748 & 5.4 & 12,416 \\
4 (3) & 11 & 19, 6,840 & 3.86 & 19 & -- & -- & -- \\
4 (3) & 12 & 19, 6,840 & 1.29 & 37 & -- & -- & -- \\
6 (5) & 7 & 13, 2,184 & 5.59 & 13 & 8,310 & 3.8 & 8,880 \\
6 (5) & 8 & 23, 12,144 & 6.22 & 31 & 29,790 & 2.5 & 31,680 \\
6 (5) & 9 & 29, 24,360 & 2.49 & 31 & 117,360 & 4.8 & 127,296 \\
6 (5) & 10 & 53, 148,824 & 3.05 & 79 & 452,032 & 3.0 & 524,160 \\
6 (5) & 11 & 67, 300,696 & 1.23 & 79 & -- & -- & -- \\
6 (5) & 12 & 107, 1,224,936 & 1.00 & 173 & -- & -- & -- \\
8 (7) & 7 & 23, 12,144 & 4.42 & 23 & 40,586 & 3.3 & 43,344 \\
8 (7) & 8 & 31, 29,760 & 1.55 & 67 & 201,480 & 6.8 & 239,904 \\
8 (7) & 9 & 67, 300,696 & 2.24 & 67 & 1,091,232 & 3.6 & 1,341,600 \\
8 (7) & 10 & 107, 1,224,936 & 1.30 & 181 & 6,927,210 & 5.7 &
7,818,096 \\
8 (7) & 11 & 173, 5,177,544 & 0.79 & 181 & -- & -- & -- \\
8 (7) & 12 & 229, 12,008,760 (scan limit) & 0.26 & 479 & -- & -- & -- \\
12 (11) & 7 & 47, 103,776 & 2.66 & 47 & 664,300 & 6.4 & 693,500 \\
12 (11) & 8 & 71, 357,840 & 0.83 & 167 & 4,650,100 & 13.0 & -- \\
12 (11) & 9 & 167, 4,657,296 & 0.99 & 167 & 30,001,152 & 6.4 &
34,026,720 \\
12 (11) & 10 & 251, 15,813,000 & 0.30 & 563 & 300,383,050 & 19.0 &
325,832,976 \\
12 (11) & 11 & 257, 16,974,336 (scan limit) & 0.03 & 563 & -- & -- &
-- \\
12 (11) & 12 & 257, 16,974,336 (scan limit) & 0.00 & 1879 & -- & -- &
-- \\
14 (13) & 7 & 59, 205,320 & 1.96 & 61 & 1,771,560 & 8.6 & 1,789,588 \\
14 (13) & 8 & 109, 1,294,920 & 0.95 & 239 & 14,172,480 & 10.9 & -- \\
14 (13) & 9 & 233, 12,649,104 & 0.72 & 239 & 95,087,738 & 7.5 &
116,068,656 \\
14 (13) & 10 & 257, 16,974,336 (scan limit) & 0.07 & 859 & 1,213,477,190
& 71.5 & 1,326,627,120 \\
14 (13) & 11 & 257, 16,974,336 (scan limit) & 0.01 & 859 & -- & -- &
-- \\
14 (13) & 12 & 257, 16,974,336 (scan limit) & 0.00 & 3089 & -- & -- &
-- \\
\end{longtable}

Scan limits: p = 3: q ≤ 157, p = 5: q ≤ 157, p = 7: q ≤ 229, p = 11: q ≤
257, p = 13: q ≤ 257.

At odd D, in the 13 cells not cut by the scan limit, the largest q
reached equals the limit of Theorem 3.2 in 8 cells ((4,7), (4,9),
(4,11), (6,7), (8,7), (8,9), (12,7), (12,9)), is the next prime below it
in 3 ((6,9): 29 against 31; (14,7): 59 against 61; (14,9): 233 against
239), and is farther below in 2 ((6,11): 67 against 79; (8,11): 173
against 181). So at odd D the vertex ψ(i) is usually what stops the
family, as far as the scan goes. At even D = 2κ the limit is the same as
at D = 2κ + 1, and the generic vertices, at distance about log\_p n,
stop the family first.

\end{document}
